\documentclass[12pt]{article}
\usepackage[utf8]{inputenc}
\usepackage[T1]{fontenc}
\usepackage{hyperref}
\title{T\`ay Language Resources for K-12 Curriculum Development}
\author{Compiled by Claude Sonnet 5 (claude-sonnet-5) for Andr\'as Kornai}
\date{}
\begin{document}
\maketitle

\section{Introduction}
T\`ay (ISO 639-3 tyz) is a Kra-Dai (Tai) language of northern Vietnam with roughly 1.6 million speakers. The searches done for this report found only general overview material and no T\`ay textbook, primer or school programme. Searches in Vietnamese (for example ``t\`u \d{\~i}\'en ti\'ang T\`ay'') would be the next step.

\section{Primary Educational Resources}
No T\`ay primer, textbook or school curriculum turned up in the searches. A 2018 article on the status and role of the T\`ay language in northeast Vietnam, which argues for stronger community attitudes toward the language, is listed under Linguistic and Reference Materials. A 2026 arXiv paper on a Cham, Khmer and Tay-Nung corpus notes that Vietnamese overwhelmingly dominates as the language of instruction.

\section{Linguistic and Reference Materials}
\begin{itemize}
\item Omniglot overview of T\`ay and its writing systems (chu N\^om T\`ay and the Latin-based script of 1961): \url{https://omniglot.com/writing/tay.htm}
\item Wikipedia entry on the T\`ay language: \url{https://wikipedia.com/wiki/Tay_language}
\item Multilingual corpus paper covering Cham, Khmer and Tay-Nung (2026, arXiv); it mentions 1970s and 1980s grammars and dictionaries without naming them: \url{https://arxiv.org/pdf/2607.08362}
\item Unicode proposal documents on the T\`ay script: \url{https://www.unicode.org/L2/L2022/22208-yo-lai-tay-script.pdf} and \url{https://www.unicode.org/L2/L2022/22152-yo-lai-tay-script.pdf}
\item International Journal of Language Teaching and Education, July 2018, article on the T\`ay language in northeast Vietnam: \url{https://doaj.org/article/e90a9e1e91c448058e1c9e8473f5601b}
\end{itemize}

\section{Religious and Community Resources}
Nothing was found in the searches.

\section{Digital and Audio Resources}
No T\`ay audio, video or software resources were found. The Unicode proposals above are relevant to digital text support for the traditional script.

\section{Curriculum Development Context}
The searches indicate that T\`ay has an official Latin orthography and some earlier documentation, but they found no evidence of T\`ay-medium schooling or published classroom materials. A curriculum would probably need to start from new materials. Searching Vietnamese-language sources is the recommended next step.
\end{document}